\documentclass[11pt]{article}
\usepackage{amsmath, amssymb, amsfonts}
\usepackage{booktabs}
\usepackage{geometry}
\usepackage{hyperref}
\usepackage{natbib}
\geometry{margin=1in}

\title{Mathematical Model for an Interactive Physarum Slime Mold Simulator}
\author{Generated for the Streamlit Physarum Simulator}
\date{\today}

\begin{document}
\maketitle

\section{Introduction}
This document defines the mathematical model implemented by the interactive
Streamlit simulator for \emph{Physarum polycephalum}. The simulator combines a
local multi-agent model with trail fields, food attractants, energy dynamics,
growth, division, dormancy, death, and a graph-based adaptive transport-network
solver. This simulator is an idealized computational model. It is inspired by
Physarum behavior and published mathematical models, but it is not a calibrated
biological model.

\section{Biological and Computational Motivation}
The maze-solving experiments of \citet{Nakagaki2000Maze} motivate the idea that
Physarum can restructure its body and transport tubes to connect nutrient
locations efficiently. The adaptive conductance network model of
\citet{Tero2007Adaptive} and the network-design rules of
\citet{Tero2010Rules} motivate the pressure, flow, and tube-conductance layer.
The multi-agent sensing and trail-field approach follows the computational
spirit of the Jones models \citep{Jones2010Pattern,Jones2011Influences}. The
shortest-path overlay is included because standard Physarum models are known to
have shortest-path convergence properties in mathematical settings
\citep{Bonifaci2013Shortest}. Finally, the food quality and nutritional
decision variables are idealized abstractions inspired by the nutritional
choice behavior described by \citet{Dussutour2010Nutrition}.

\section{Domain, Time, and State Space}
Time is discrete:
\begin{equation}
  t_n = n\Delta t,\qquad n=0,1,2,\ldots,
\end{equation}
with default step size $\Delta t=1$. The discrete computational domain is
\begin{equation}
  \Omega = \{0,1,\ldots,W-1\}\times \{0,1,\ldots,H-1\}.
\end{equation}
Agent positions are continuous points in the embedded rectangle
\begin{equation}
  \bar{\Omega}=[0,W)\times[0,H).
\end{equation}
Fields are stored on the grid $\Omega$. For a grid field $F$ and a continuous
point $z=(x,y)\in\bar{\Omega}$, the simulator uses bilinear interpolation,
written $\mathcal{I}[F](z)$. If $x_0=\lfloor x\rfloor$, $y_0=\lfloor y\rfloor$,
$\delta_x=x-x_0$, and $\delta_y=y-y_0$, then
\begin{align}
  \mathcal{I}[F](x,y)
  &=
  (1-\delta_x)(1-\delta_y)F_{y_0,x_0}
  +\delta_x(1-\delta_y)F_{y_0,x_0+1} \nonumber\\
  &\quad +(1-\delta_x)\delta_y F_{y_0+1,x_0}
  +\delta_x\delta_y F_{y_0+1,x_0+1},
\end{align}
with boundary indices either wrapped or clipped according to the boundary mode.

\section{Agent State Variables}
There are $N(t)$ agents at time $t$. Agent $i$ has state
\begin{equation}
  X_i(t)=\left(p_i(t),\theta_i(t),E_i(t),m_i(t),a_i(t),\mu_i(t)\right).
\end{equation}
The variables are:
\begin{itemize}
  \item $p_i(t)=(x_i(t),y_i(t))\in\bar{\Omega}$, the continuous position.
  \item $\theta_i(t)\in(-\pi,\pi]$, the heading angle.
  \item $E_i(t)\in[0,E_{\max}]$, the internal energy.
  \item $m_i(t)>0$, the biomass or mass.
  \item $a_i(t)\in\{0,1\}$, the alive indicator.
  \item $\mu_i(t)\in\{\mathrm{search},\mathrm{exploit},\mathrm{dormant}\}$,
  the behavior mode.
\end{itemize}
The implementation stores positions as an $N\times 2$ array and headings,
energies, masses, alive indicators, and mode labels as one-dimensional arrays.

\section{Food Source Variables}
There are $M(t)$ food sources. Food source $k$ has state
\begin{equation}
  F_k(t)=\left(f_k,C_k(t),q_k,\sigma_{A,k},r_{F,k}\right),
\end{equation}
where $f_k=(u_k,v_k)\in\bar{\Omega}$ is the location, $C_k(t)\ge 0$ is remaining
abstract calories, $q_k>0$ is a quality or attractiveness multiplier, and
$\sigma_{A,k}>0$ is the attractant spread scale. The parameter $r_{F,k}>0$ is
an interactive food size used for visualization and for widening the local
consumption neighborhood. A food source is active when
\begin{equation}
  C_k(t)>C_{\min}.
\end{equation}

\section{Environment Fields}
The simulator maintains four grid fields:
\begin{align}
  T_t(x,y)&\ge 0 && \text{trail or chemoattractant deposited by agents},\\
  A_t(x,y)&\ge 0 && \text{food attractant},\\
  O(x,y)&\in\{0,1\} && \text{obstacle mask, with 1 meaning blocked},\\
  R_t(x,y)&\ge 0 && \text{optional repellent or hazard field}.
\end{align}
The combined sensory field is
\begin{equation}
  S_t(x,y)=w_TT_t(x,y)+w_AA_t(x,y)-w_OO(x,y)-w_RR_t(x,y).
\end{equation}
The default obstacle weight is large, $w_O=10^3$, so obstacle cells are strongly
repellent.

\section{Gaussian Diffusion and Field Sampling}
The ideal discrete Gaussian kernel is
\begin{equation}
  G_\sigma(r)=\frac{1}{Z_\sigma}\exp\left(-\frac{\|r\|_2^2}{2\sigma^2}\right),
  \qquad r\in\mathbb{Z}^2,
\end{equation}
where $Z_\sigma$ normalizes the kernel so that
\begin{equation}
  \sum_{r\in\mathbb{Z}^2}G_\sigma(r)=1.
\end{equation}
In code, convolution by $G_\sigma$ is implemented using
\texttt{scipy.ndimage.gaussian\_filter}. The trail update is
\begin{equation}
  T_{t+\Delta t}=(1-\lambda_T\Delta t)
  \left[G_{\sigma_T}*(T_t+D_t)\right],
\end{equation}
followed by nonnegative clipping:
\begin{equation}
  T_{t+\Delta t}(x,y)\leftarrow \max\{0,T_{t+\Delta t}(x,y)\}.
\end{equation}

\section{Trail Deposition}
Each alive, non-dormant agent deposits trail at its current position. The
deposition amount is
\begin{equation}
  q_i^T(t)=q_0\,m_i(t)\,\phi_E(E_i(t))\,\phi_\mu(\mu_i(t)),
\end{equation}
where
\begin{equation}
  \phi_E(E)=\operatorname{clip}\left(\frac{E}{E_{\max}},0,1\right)
\end{equation}
and
\begin{equation}
  \phi_\mu(\mu)=
  \begin{cases}
  1, & \mu=\mathrm{search},\\
  \rho_{\mathrm{exploit}}, & \mu=\mathrm{exploit},\\
  0, & \mu=\mathrm{dormant}.
  \end{cases}
\end{equation}
The default is $\rho_{\mathrm{exploit}}=1.2$. The deposition field is
\begin{equation}
  D_t(x,y)=\sum_{i:a_i(t)=1}q_i^T(t)K_h((x,y)-p_i(t)).
\end{equation}
The implementation uses bilinear splatting. If an agent lies at $(x,y)$ with
fractional offsets $\delta_x$ and $\delta_y$, its amount is split among the four
neighboring grid cells with weights
\begin{equation}
  (1-\delta_x)(1-\delta_y),\quad
  \delta_x(1-\delta_y),\quad
  (1-\delta_x)\delta_y,\quad
  \delta_x\delta_y.
\end{equation}

\section{Food Attractant Field}
At each step, active food sources define the attractant
\begin{equation}
  A_t(x,y)=\sum_{k=1}^{M(t)}
  q_k h(C_k(t))
  \exp\left(-\frac{\|(x,y)-f_k\|_2^2}{2\sigma_{A,k}^2}\right).
\end{equation}
The default calorie response is saturating:
\begin{equation}
  h(C)=\frac{C}{C+C_{1/2}+\varepsilon}.
\end{equation}
This prevents very large food sources from dominating the entire world. A
simpler optional conceptual transform is $h(C)=C/C_{\mathrm{scale}}$, but the
implementation uses the saturating form by default.

\section{Agent Sensing}
For heading $\theta$, define
\begin{equation}
  u(\theta)=
  \begin{bmatrix}
  \cos\theta\\
  \sin\theta
  \end{bmatrix}.
\end{equation}
The front, left, and right sensor coordinates are
\begin{align}
  z_i^0(t)&=p_i(t)+d_su(\theta_i(t)),\\
  z_i^L(t)&=p_i(t)+d_su(\theta_i(t)+\alpha),\\
  z_i^R(t)&=p_i(t)+d_su(\theta_i(t)-\alpha).
\end{align}
Their sampled values are
\begin{align}
  S_i^0(t)&=\mathcal{I}[S_t](z_i^0(t)),\\
  S_i^L(t)&=\mathcal{I}[S_t](z_i^L(t)),\\
  S_i^R(t)&=\mathcal{I}[S_t](z_i^R(t)).
\end{align}
If a sensor lies outside the world in bounce mode, its value is set to a strong
repellent value. In wrap mode, sensor coordinates are taken modulo $(W,H)$.

\section{Turning and Motion}
Let
\begin{equation}
  S_i^{\max}(t)=\max\{S_i^0(t),S_i^L(t),S_i^R(t)\}.
\end{equation}
Let $\delta_i(t)\in\{0,+1,-1\}$ select the largest sensor, where $0$ denotes
front, $+1$ left, and $-1$ right. The deterministic turn is
\begin{equation}
  \Delta\theta_i^{\mathrm{det}}(t)=
  \begin{cases}
  +\beta, & \delta_i(t)=+1,\\
  -\beta, & \delta_i(t)=-1,\\
  0, & \delta_i(t)=0.
  \end{cases}
\end{equation}
If the agent is searching or the local maximum signal is weak, the model adds
random exploration $\xi_i(t)\sim\mathcal{N}(0,\sigma_\theta^2)$. The heading is
\begin{equation}
  \theta_i(t+\Delta t)=
  \operatorname{wrap}_{(-\pi,\pi]}
  \left(\theta_i(t)+\Delta\theta_i^{\mathrm{det}}(t)+\chi_i(t)\xi_i(t)\right),
\end{equation}
where $\chi_i(t)=1$ for search mode or $S_i^{\max}(t)<\tau_S$ and is zero
otherwise.

The energy speed factor is
\begin{equation}
  \psi_E(E)=\operatorname{clip}
  \left(
  \frac{v_{\min}}{v_0}
  +\left(1-\frac{v_{\min}}{v_0}\right)\frac{E}{E_{\max}},
  \frac{v_{\min}}{v_0},1
  \right),
\end{equation}
and the speed is
\begin{equation}
  v_i(t)=v_0\psi_E(E_i(t))\psi_\mu(\mu_i(t)),
\end{equation}
where
\begin{equation}
  \psi_\mu(\mu)=
  \begin{cases}
  1, & \mu=\mathrm{search},\\
  v_{\mathrm{exploit}}, & \mu=\mathrm{exploit},\\
  0, & \mu=\mathrm{dormant}.
  \end{cases}
\end{equation}
The proposed movement is
\begin{equation}
  \tilde p_i(t+\Delta t)=p_i(t)+v_i(t)\Delta t\,u(\theta_i(t+\Delta t)).
\end{equation}
Boundary handling is either wrapping modulo $(W,H)$ or reflection with heading
updates. If the proposed point enters an obstacle cell, movement is rejected and
the heading is rotated:
\begin{equation}
  \theta_i(t+\Delta t)\leftarrow
  \operatorname{wrap}\left(\theta_i(t+\Delta t)+\pi+\eta_i\right),
  \qquad \eta_i\sim U(-\pi/4,\pi/4).
\end{equation}

\section{Search versus Exploit Decision}
The logistic sigmoid is
\begin{equation}
  \sigma(x)=\frac{1}{1+\exp(-x)}.
\end{equation}
The energy sufficiency score is
\begin{equation}
  P_i^E(t)=
  \sigma\left(k_E\left[\frac{E_i(t)}{E_{\max}}-\theta_E\right]\right),
\end{equation}
and the signal exploitation score is
\begin{equation}
  P_i^S(t)=
  \sigma\left(k_S\left[S_i^{\max}(t)-\tau_S\right]\right).
\end{equation}
The probability of broad searching is
\begin{equation}
  P_i^{\mathrm{search}}(t)=P_i^E(t)\left[1-P_i^S(t)\right].
\end{equation}
The mode update is
\begin{equation}
  \mu_i(t+\Delta t)=
  \begin{cases}
  \mathrm{dormant}, &
  \frac{E_i(t)}{E_{\max}}<\theta_{\mathrm{dormant}}
  \text{ and } S_i^{\max}(t)<\tau_S,\\
  \mathrm{search}, & U_i(t)<P_i^{\mathrm{search}}(t),\\
  \mathrm{exploit}, & \text{otherwise},
  \end{cases}
\end{equation}
where $U_i(t)\sim U(0,1)$.
When the optional ``continue searching without food'' switch is enabled and no
active food sources remain, non-dormant agents are forced into search mode. This
keeps the population exploring depleted environments instead of merely following
stale trail.

\section{Food Consumption}
Agent $i$ can eat food source $k$ if
\begin{equation}
  \|p_i(t)-f_k\|_2\le r_{\mathrm{eat}}+r_{F,k}.
\end{equation}
The raw demand is
\begin{equation}
  d_{ik}(t)=b_{\max}\Delta t\,
  \mathbf{1}[a_i(t)=1]\,
  \mathbf{1}[\mu_i(t)\ne\mathrm{dormant}]\,
  \mathbf{1}[\|p_i(t)-f_k\|_2\le r_{\mathrm{eat}}+r_{F,k}].
\end{equation}
If multiple agents request the same food source, calories are allocated
proportionally:
\begin{equation}
  b_{ik}(t)=
  \begin{cases}
  d_{ik}(t)\min\left(1,\frac{C_k(t)}{\sum_j d_{jk}(t)+\varepsilon}\right),
  & \sum_j d_{jk}(t)>0,\\
  0, & \sum_j d_{jk}(t)=0.
  \end{cases}
\end{equation}
Total eaten from food source $k$ is
\begin{equation}
  B_k(t)=\sum_i b_{ik}(t),
\end{equation}
and food calories update by
\begin{equation}
  C_k(t+\Delta t)=\max(0,C_k(t)-B_k(t)).
\end{equation}
The energy gained by agent $i$ is
\begin{equation}
  G_i^{\mathrm{food}}(t)=\eta_{\mathrm{food}}\sum_{k=1}^{M(t)}q_kb_{ik}(t).
\end{equation}
The direct implementation checks distances between agents and each active food
source. This is $O(NM)$ per consumption update, which is acceptable for the
moderate number of food sources targeted by the interactive app.

\section{Energy Budget}
The metabolic cost is
\begin{equation}
  L_i^{\mathrm{base}}(t)=c_{\mathrm{base}}m_i(t)\Delta t.
\end{equation}
The movement distance is
\begin{equation}
  \Delta s_i(t)=\|p_i(t+\Delta t)-p_i(t)\|_2,
\end{equation}
and movement cost is
\begin{equation}
  L_i^{\mathrm{move}}(t)=c_{\mathrm{move}}m_i(t)\Delta s_i(t).
\end{equation}
Trail-deposition cost is
\begin{equation}
  L_i^{\mathrm{trail}}(t)=c_{\mathrm{trail}}q_i^T(t).
\end{equation}
Dormant cost is
\begin{equation}
  L_i^{\mathrm{dormant}}(t)=c_{\mathrm{dormant}}m_i(t)\Delta t.
\end{equation}
Total cost is
\begin{equation}
  L_i(t)=L_i^{\mathrm{base}}+L_i^{\mathrm{move}}+
  L_i^{\mathrm{trail}}+\mathbf{1}[\mu_i=\mathrm{dormant}]L_i^{\mathrm{dormant}}.
\end{equation}
Energy before growth is
\begin{equation}
  \hat E_i(t+\Delta t)=
  E_i(t)+G_i^{\mathrm{food}}(t)-L_i(t),
\end{equation}
followed by clipping:
\begin{equation}
  \hat E_i(t+\Delta t)\leftarrow
  \operatorname{clip}\left(\hat E_i(t+\Delta t),0,E_{\max}\right).
\end{equation}

\section{Growth and Division}
Growth is possible only when energy is above a threshold. With
$[x]_+=\max(x,0)$, the specific growth rate is
\begin{equation}
  r_i^g(t)=r_g
  \left[\frac{\hat E_i(t+\Delta t)}{E_{\max}}-\theta_g\right]_+.
\end{equation}
Mass increase is
\begin{equation}
  \Delta m_i(t)=r_i^g(t)m_i(t)\Delta t,
\end{equation}
and growth energy cost is
\begin{equation}
  L_i^{\mathrm{growth}}(t)=c_{\mathrm{growth}}\Delta m_i(t).
\end{equation}
The pre-division mass and energy are
\begin{align}
  \tilde m_i(t+\Delta t)&=m_i(t)+\Delta m_i(t),\\
  \tilde E_i(t+\Delta t)&=
  \operatorname{clip}\left(\hat E_i(t+\Delta t)-L_i^{\mathrm{growth}}(t),0,E_{\max}\right).
\end{align}
If $\tilde E_i(t+\Delta t)>E_{\mathrm{split}}$ and $N(t)<N_{\max}$, then the
agent divides. The child position is
\begin{equation}
  p_j(t+\Delta t)=p_i(t+\Delta t)+\epsilon_p,\qquad
  \epsilon_p\sim\mathcal{N}(0,\sigma_{\mathrm{split}}^2I_2).
\end{equation}
The child heading is
\begin{equation}
  \theta_j(t+\Delta t)=
  \operatorname{wrap}\left(\theta_i(t+\Delta t)+U(-\pi,\pi)\right).
\end{equation}
Energy and mass are split evenly:
\begin{equation}
  E_i(t+\Delta t)=E_j(t+\Delta t)=\frac{1}{2}\tilde E_i(t+\Delta t),
\end{equation}
\begin{equation}
  m_i(t+\Delta t)=m_j(t+\Delta t)=\frac{1}{2}\tilde m_i(t+\Delta t).
\end{equation}
If more agents are eligible than the remaining capacity, eligible agents are
randomly selected until $N_{\max}$ is reached.

\section{Death and Dormancy}
If
\begin{equation}
  E_i(t+\Delta t)\le E_{\mathrm{death}},
\end{equation}
the user-selected policy is applied. Under the remove policy,
\begin{equation}
  a_i(t+\Delta t)=0.
\end{equation}
Under the dormant policy,
\begin{equation}
  \mu_i(t+\Delta t)=\mathrm{dormant},\qquad E_i(t+\Delta t)=0.
\end{equation}
The default is removal for visual clarity.

\section{Adaptive Physarum Transport Network}
At selected intervals, the simulator builds a graph
\begin{equation}
  G_t=(V_t,E_t).
\end{equation}
Nodes include a nest node, active food nodes, optional high-trail landmarks, and
optional support nodes. The nest is the energy-weighted colony center:
\begin{equation}
  c_t=\frac{\sum_i a_i(t)E_i(t)p_i(t)}{\sum_i a_i(t)E_i(t)+\varepsilon}.
\end{equation}
If all agents are dead, the previous colony center or geometric center is used.
Edges are constructed by $k$-nearest-neighbor links, with optional Delaunay-like
links when available. Obstacle-crossing edges can be rejected by line sampling.
Each edge $e=(u,v)$ has length
\begin{equation}
  \ell_e=\|x_u-x_v\|_2,
\end{equation}
conductance $D_e(t)>0$, signed flow $Q_e(t)$, and path cost $c_e(t)$.

\section{Pressure, Flow, and Conductance Dynamics}
Edges are oriented arbitrarily. The incidence matrix is
\begin{equation}
  B_{ve}=
  \begin{cases}
  +1, & v=\mathrm{head}(e),\\
  -1, & v=\mathrm{tail}(e),\\
  0, & \text{otherwise}.
  \end{cases}
\end{equation}
The hydraulic coefficient is
\begin{equation}
  g_e(t)=\frac{D_e(t)}{\ell_e+\varepsilon},
\end{equation}
and
\begin{equation}
  G_D(t)=\operatorname{diag}(g_e(t)).
\end{equation}
The weighted graph Laplacian is
\begin{equation}
  L_D(t)=BG_D(t)B^\top.
\end{equation}
The source-sink vector $b(t)$ assigns food nodes positive injection:
\begin{equation}
  b_k(t)=I_0\frac{C_k(t)}{\sum_{\ell\in\mathcal{F}_t}C_\ell(t)+\varepsilon},
\end{equation}
and the nest is the sink:
\begin{equation}
  b_0(t)=-\sum_{k\in\mathcal{F}_t}b_k(t).
\end{equation}
Intermediate nodes have $b_v(t)=0$, so $\sum_vb_v(t)=0$.

Pressures solve
\begin{equation}
  L_D(t)P(t)=b(t).
\end{equation}
Because $L_D$ is singular, the gauge $P_0(t)=0$ is fixed by removing the nest row
and column. For edge $e=(u,v)$, pressure drop and flow are
\begin{align}
  \Delta P_e(t)&=P_u(t)-P_v(t),\\
  Q_e(t)&=\frac{D_e(t)}{\ell_e+\varepsilon}\left[P_u(t)-P_v(t)\right].
\end{align}
Conductance adapts by a discrete Physarum-like tube equation:
\begin{equation}
  \frac{dD_e}{dt}=\alpha_D f(|Q_e|)-\mu_DD_e,
\end{equation}
where
\begin{equation}
  f(q)=\frac{q^\gamma}{q^\gamma+q_0^\gamma+\varepsilon}.
\end{equation}
The implemented update is
\begin{equation}
  D_e(t+\Delta t)=
  \max\left(
  D_{\min},
  D_e(t)+\Delta t\left[
  \alpha_Df(|Q_e(t)|)-\mu_DD_e(t)
  \right]
  \right).
\end{equation}
This adaptive conductance equation is the primary optimization mechanism:
high-flow edges are reinforced and low-flow edges decay.

\section{Shortest Nutrient Path Overlay}
After conductance adaptation, the visualization computes Dijkstra shortest
paths only as an analysis layer. The path cost is
\begin{equation}
  c_e^{\mathrm{path}}(t)=\frac{\ell_e}{(D_e(t)+\varepsilon)^\eta}.
\end{equation}
The primary food node is
\begin{equation}
  k^\star=\arg\max_k C_k(t),
\end{equation}
and the overlay path is
\begin{equation}
  \pi^\star(t)=\operatorname{Dijkstra}(G_t,v_0,v_{k^\star},c^{\mathrm{path}}).
\end{equation}
The geometric path length and weighted path cost are
\begin{align}
  L_{\mathrm{path}}(t)&=\sum_{e\in\pi^\star(t)}\ell_e,\\
  C_{\mathrm{path}}(t)&=\sum_{e\in\pi^\star(t)}c_e^{\mathrm{path}}(t).
\end{align}

\section{Coupling Between Network and Agent Field}
The agent and network models are coupled in two directions. First, high-trail
regions can become landmark nodes in the graph. Second, reinforced network edges
can be rasterized back into the trail field. For edge $e$, define normalized
reinforcement
\begin{equation}
  R_e(t)=\frac{|Q_e(t)|}{\max_{e'}|Q_{e'}(t)|+\varepsilon}.
\end{equation}
Grid cells near the edge receive
\begin{equation}
  T_t(x)\leftarrow T_t(x)+
  \rho_{\mathrm{net}}D_e(t)R_e(t)
  K_{\sigma_{\mathrm{edge}}}(\operatorname{dist}(x,e)).
\end{equation}
This lets flow-carrying network links bias future agent motion through the same
trail field used by local sensing.

\section{Metrics}
The simulator tracks:
\begin{align}
  N_{\mathrm{alive}}(t)&=\sum_i a_i(t),\\
  \bar E(t)&=\frac{\sum_i a_i(t)E_i(t)}{N_{\mathrm{alive}}(t)+\varepsilon},\\
  E_{\mathrm{tot}}(t)&=\sum_i a_i(t)E_i(t),\\
  M_{\mathrm{bio}}(t)&=\sum_i a_i(t)m_i(t),\\
  C_{\mathrm{food}}(t)&=\sum_k C_k(t),\\
  C_{\mathrm{consumed}}(t)&=C_{\mathrm{food}}(0)-C_{\mathrm{food}}(t),\\
  T_{\mathrm{mass}}(t)&=\sum_{x\in\Omega}T_t(x),\\
  A_{\mathrm{trail}}(t)&=\#\{x\in\Omega:T_t(x)>\tau_T\}.
\end{align}
Trail entropy is
\begin{equation}
  H_T(t)=-\sum_x\tilde T_t(x)\log(\tilde T_t(x)+\varepsilon),
  \qquad
  \tilde T_t(x)=\frac{T_t(x)}{\sum_yT_t(y)+\varepsilon}.
\end{equation}
For the network, the active edge set is
\begin{equation}
  E_t^+=\{e\in E_t:D_e(t)>D_{\mathrm{vis}}\}.
\end{equation}
Total geometric length and weighted material cost are
\begin{align}
  L_{\mathrm{net}}(t)&=\sum_{e\in E_t^+}\ell_e,\\
  C_{\mathrm{net}}(t)&=\sum_{e\in E_t}\ell_eD_e(t).
\end{align}
Average transport distance and efficiency are
\begin{align}
  \bar d_{\mathrm{net}}(t)&=
  \frac{2}{|V_t|(|V_t|-1)}\sum_{u<v}d_{uv}(t),\\
  \mathcal{E}_{\mathrm{net}}(t)&=
  \frac{2}{|V_t|(|V_t|-1)}\sum_{u<v}\frac{1}{d_{uv}(t)+\varepsilon}.
\end{align}
Disconnected pairs are ignored in the implementation rather than assigned
infinite distance. Fault tolerance is
\begin{equation}
  \mathrm{FT}(t)=
  \frac{
  \#\{e\in E_t^+:G_t^+\setminus e\text{ still connects all active food nodes to the nest}\}
  }{|E_t^+|+\varepsilon}.
\end{equation}

\section{Parameter Table}
\begin{center}
\begin{tabular}{lll}
\toprule
Category & Symbol or name & Default \\
\midrule
World & $W,H$ & $120,80$ \\
World & boundary mode & bounce \\
Agents & initial agents & $800$ \\
Agents & $N_{\max}$ & $3000$ \\
Agents & $E_{\max}$ & $100$ \\
Movement & $v_0$ & $1.0$ \\
Movement & $v_{\min}$ & $0.2$ \\
Sensing & $d_s,\alpha,\beta$ & $5.0,0.7,0.35$ \\
Sensing & $\sigma_\theta$ & $0.5$ \\
Fields & $w_T,w_A,w_O$ & $1.0,2.5,1000$ \\
Trail & $q_0,\lambda_T,\sigma_T$ & $1.0,0.04,0.8$ \\
Food & $r_{\mathrm{eat}},b_{\max},\eta_{\mathrm{food}}$ & $2.5,2.0,0.8$ \\
Food & $r_F,\sigma_A,C_{1/2}$ & $3.0,10.0,200.0$ \\
Energy & $c_{\mathrm{base}},c_{\mathrm{move}}$ & $0.03,0.05$ \\
Growth & $r_g,\theta_g,E_{\mathrm{split}}$ & $0.02,0.65,85.0$ \\
Search & $k_E,\theta_E,k_S,\tau_S$ & $8.0,0.25,5.0,0.15$ \\
Network & $I_0,\alpha_D,\mu_D$ & $2.0,1.0,0.08$ \\
Network & $\gamma,q_0,D_{\min},D_{\mathrm{init}}$ & $1.4,1.0,10^{-4},0.05$ \\
Network & $\eta,\rho_{\mathrm{net}}$ & $1.0,0.1$ \\
\bottomrule
\end{tabular}
\end{center}

\section{Algorithm Summary}
Each simulation step performs the following operations:
\begin{enumerate}
  \item Recompute the food attractant field from active food sources.
  \item Form the combined sensory field from trail, food, obstacles, and repellent.
  \item Sense front, left, and right values for each alive agent using bilinear interpolation.
  \item Update search, exploit, or dormant mode from energy and local signal.
  \item Turn, move, handle boundaries, and reject obstacle collisions.
  \item Deposit, diffuse, decay, and clip trail.
  \item Consume nearby food with proportional sharing.
  \item Apply energy gains and costs.
  \item Apply growth, division, death, or dormancy.
  \item At network intervals, build or update the graph, solve pressure and flow, adapt conductances, compute Dijkstra overlays, and rasterize network reinforcement into the trail field.
  \item Record scalar metrics.
\end{enumerate}

\section{Model Limitations}
The model is intentionally idealized. It does not resolve detailed cytoplasmic
streaming, tube wall mechanics, intracellular biochemical oscillations,
multi-nutrient stoichiometry, humidity, light exposure, microbial competition,
or laboratory-scale calibration. Food calories are abstract simulation units.
Agent division is a computational population-control mechanism rather than a
complete biological account of Physarum growth. The adaptive network solver is
coupled to the agent field but remains a graph abstraction, not a physical tube
mesh. Dijkstra paths are diagnostic overlays and should not be interpreted as
the mechanism that creates the adaptive network.

\section{References}
\bibliographystyle{plainnat}
\bibliography{references}

\end{document}
